\section{\texorpdfstring{The sixth sum at \(N/4\)}{The sixth sum at N/4}}\label{sec:root-a}

Put \(L=\log N\) and fix \(0<\sigma\le10^{-10}\) for the common
error allocation. No estimate in the sixth-sum construction depends on \(\sigma\). Fix
\begin{equation}\label{eq:ra-parameters}
\begin{gathered}
 \alpha=\frac{100}{1149},\qquad t_0=\frac{57}{224},\qquad
 v_{\mathrm{RA}}=\frac{5227}{10000},\qquad \eta_{\mathrm m}=\varepsilon=\frac1{10000},\\
 J=[t_0,1/3),\quad J_t=[t_0,1/3-\eta_{\mathrm m}),\quad
 J_{\mathrm{tail}}=[1/3-\eta_{\mathrm m},1/3),\\
 I_L=[2,N/4],\qquad I_H=(N/4,(1-\varepsilon)N),\qquad z_s=N^\alpha.
\end{gathered}
\end{equation}
Write \(p\in N^K\) for \(p\in\mathbb P\) with \(\log p/L\in K\).
Paper II, (1.7), gives the original sum \cite[(1.7)]{PaperII}:
\begin{equation}\label{eq:ra-original-s6}
 S_6=\sum_{\substack{p\in N^J\\(p,N)=1}}
 \ \sum_{\substack{m\ge1,\ e=N-pm\in\mathbb P\\2\le e<(1-\varepsilon)N}}
 \mathbf1_{P^-(m)\ge z_s}\mathbf1_{(m,Np)=1}.
\end{equation}
Each prime mark remains a separate occurrence.

\subsection{Three upper weights}

For squarefree \(d=\ell_1\cdots\ell_k\), order its primes decreasingly.
The ordinary upper Rosser weight is
\begin{equation}\label{eq:ra-rosser-weight}
 \mu_d^+(D)=\mu(d)
 \prod_{\substack{1\le j\le k\\j\ \mathrm{odd}}}
 \mathbf1_{\ell_1\cdots\ell_{j-1}\ell_j^3<D},
 \qquad \mu_1^+(D)=1.
\end{equation}
The weight is zero outside its stated squarefree support.
Its divisor sum is a nonnegative upper sieve weight
\cite[Lemma 1, (3.2), (3.6)]{Iwaniec1980Rosser}
\cite[Lemmas 2--3, p.~313]{Iwaniec1980Remainder}.

For \(K\subseteq J\) and \(p\in N^K\), define
\begin{equation}\label{eq:ra-coefficient}
 c_{R,K}(pd)=\mu(d)\mathbf1_{pd<R}
 \prod_{\substack{1\le j\le k\\j\ \mathrm{odd}}}
 \mathbf1_{p\ell_1\cdots\ell_{j-1}\ell_j^3<R}.
\end{equation}
Require \(\ell_j<z_s\) and \((pd,2N)=1\).
Set the coefficient to zero otherwise.
Since \(t_0>\alpha\), the mark is unique.
Thus \(|c_{R,K}(q)|\le1\).
On supported moduli with \(p<R\), the gates give \(d<R/p\).
For \(d=1\) the gate product is empty, and \(p<R\) suffices.
Every marked prime has \(p<N^{1/3}\).
Each sieve level below is \(>N^{1/3}\) for large \(N\).
Hence
\begin{equation}\label{eq:ra-weight-equality}
 c_{R,K}(pd)=\mu_d^+(R/p),\qquad
 \varphi(pd)=(p-1)\varphi(d).
\end{equation}
The equality is claimed on supported moduli only.
Off that support the coefficient is zero.

Choose fixed logarithmic exponents. Put
\begin{equation}\label{eq:ra-factories}
\begin{aligned}
 R_L&=N^{1/2}L^{-B_L},&
 R_{\mathrm{tail}}&=N^{1/2}L^{-B_{\mathrm{tail}}},\\
 Q&=N^{v_{\mathrm{RA}}}L^{-B},&Q_0&=N^{1/2}L^{-B_0},\\
 b_{\mathrm{all}}&=c_{R_L,J},&
 b_{\mathrm{tail}}&=c_{R_{\mathrm{tail}},J_{\mathrm{tail}}},\\
 c_{\mathrm{trim}}&=c_{Q,J_t},&
 c_{\mathrm{lo}}(q)&=c_{\mathrm{trim}}(q)\mathbf1_{q\le Q_0},\\
 \Gamma(q)&=c_{\mathrm{trim}}(q)\mathbf1_{q>Q_0}.
\end{aligned}
\end{equation}
The last two arrays keep the gates at \(Q\).

First omit \((m,Np)=1\) in \eqref{eq:ra-original-s6}.
Sieve the larger sequence by primes \(\ell\nmid2N\), \(\ell<z_s\).
Use the three weights on
\begin{equation}\label{eq:ra-three-pieces}
 I_L\times J,\qquad I_H\times J_{\mathrm{tail}},\qquad I_H\times J_t.
\end{equation}
Their divisor sum is the finite weighted count
\begin{equation}\label{eq:ra-mix-definition}
\begin{aligned}
 A_{\mathrm{mix}}={}&
 \sum_q b_{\mathrm{all}}(q)
       \sum_{\substack{e\in I_L\cap\mathbb P\\e\equiv N\ (q)}}1\\
 &+\sum_q\{b_{\mathrm{tail}}(q)+c_{\mathrm{trim}}(q)\}
       \sum_{\substack{e\in I_H\cap\mathbb P\\e\equiv N\ (q)}}1.
\end{aligned}
\end{equation}
The upper-weight comparison
\begin{equation}\label{eq:ra-positive-comparison}
 S_6\le A_{\mathrm{mix}},\qquad
 Z_6:=A_{\mathrm{mix}}-S_6\ge0,
\end{equation}
is proved in Lemma~\ref{lem:ra-upper-sieve}(a) below.
Positivity is used before expanding the divisor sums.

\subsection{The exact discrepancy identity}

For a bounded interval \(I\subset[2,\infty)\) and a finitely supported
modulus array \(a(q)\), put \(\operatorname{Li}(I)=\int_I dt/\log t\)
and define
\begin{equation}\label{eq:ra-discrepancies}
\begin{aligned}
 \Pi_I(q)&=\sum_{e\in I\cap\mathbb P}\mathbf1_{e\equiv N\pmod q},\\
 B_I(q)&=\frac1{\varphi(q)}
             \sum_{e\in I\cap\mathbb P}\mathbf1_{(e,q)=1},\\
 D_I(a)&=\sum_q a(q)
       \left\{\Pi_I(q)-\frac{\operatorname{Li}(I)}{\varphi(q)}\right\},\\
 P_I(a)&=\sum_q a(q)\{\Pi_I(q)-B_I(q)\},\qquad
 C_I(a)=D_I(a)-P_I(a).
\end{aligned}
\end{equation}
Also put
\begin{equation}\label{eq:ra-main-mass}
\begin{aligned}
 M_{\mathrm{mix}}={}&\operatorname{Li}(I_L)
              \sum_q\frac{b_{\mathrm{all}}(q)}{\varphi(q)}\\
 &+\operatorname{Li}(I_H)
              \sum_q\frac{b_{\mathrm{tail}}(q)+c_{\mathrm{trim}}(q)}{\varphi(q)}.
\end{aligned}
\end{equation}
Finite expansion gives
\begin{equation}\label{eq:ra-b9}
\begin{aligned}
 A_{\mathrm{mix}}={}&M_{\mathrm{mix}}+D_{I_L}(b_{\mathrm{all}})
     +D_{I_H}(b_{\mathrm{tail}})\\
 &+C_{I_H}(c_{\mathrm{trim}})+P_{I_H}(c_{\mathrm{lo}})
     +P_{I_H}(\Gamma).
\end{aligned}
\end{equation}
Since \(c_{\mathrm{trim}}=c_{\mathrm{lo}}+\Gamma\),
\begin{equation}\label{eq:ra-cancellation}
 C_{I_H}(c_{\mathrm{trim}})+P_{I_H}(c_{\mathrm{lo}})
       =D_{I_H}(c_{\mathrm{lo}})+C_{I_H}(\Gamma).
\end{equation}
Therefore
\begin{equation}\label{eq:ra-b9-f1}
\begin{aligned}
 A_{\mathrm{mix}}={}&M_{\mathrm{mix}}+D_{I_L}(b_{\mathrm{all}})
     +D_{I_H}(b_{\mathrm{tail}})+D_{I_H}(c_{\mathrm{lo}})\\
 &+C_{I_H}(\Gamma)+P_{I_H}(\Gamma).
\end{aligned}
\end{equation}
No low natural-coprime term is estimated separately.

For large \(N\), each prime in \(I_H\) exceeds \(Q\).
Hence
\begin{equation}\label{eq:ra-high-conversion}
 C_{I_H}(\Gamma)=
 \{\pi(I_H)-\operatorname{Li}(I_H)\}
          \sum_q\frac{\Gamma(q)}{\varphi(q)}.
\end{equation}
The elementary identity
\begin{equation}\label{eq:ra-totient-mass}
 \frac q{\varphi(q)}=\sum_{d\mid q}\frac{\mu(d)^2}{\varphi(d)},
 \qquad
 \sum_{q\le Q}\frac1{\varphi(q)}\ll\log(2Q)
\end{equation}
follows by multiplicativity and a convergent Euler product.
Thus Paper I, (2.21), bounds \eqref{eq:ra-high-conversion}
by \(O_A(NL^{-A})\) for each fixed \(A\)
\cite[(2.21)]{PaperI}.

For the three \(D\)-terms we use the maximal-endpoint prime form of
Bombieri--Vinogradov.
Here it is conditional on the prime Siegel--Walfisz input
\cite[(2.21)]{PaperI}. Define the prime progression error by
\begin{equation}\label{eq:ra-prime-error}
 \Delta_\pi(y;q,a)=\pi(y;q,a)-\frac{\operatorname{Li}(y)}{\varphi(q)}.
\end{equation}
Fact~\ref{fact:prime-estimates-maximal} gives, for each fixed \(S>0\),
\begin{equation}\label{eq:ra-prime-bv}
 \sum_{q\le X^{1/2}(\log X)^{-(S+60)}}
 \max_{(a,q)=1}\sup_{0\le y\le X}
 |\Delta_\pi(y;q,a)|\ll_S X(\log X)^{-S}.
\end{equation}
Here \(\operatorname{Li}(y)=0\) for \(0\le y<2\).
The level sits at \(X\), not at the variable endpoint \(y\),
so intervals with endpoints chosen separately for each
modulus cost at most twice this bound. This is the maximal prime statement.
We do not substitute the fixed-endpoint von Mangoldt statement of
Paper I, (2.20). The source and the fixed \(60\) logarithmic buffer
are those of Fact~\ref{fact:prime-estimates-maximal}.

Set \(X=N\), \(S=A+10\), and choose the half-level buffers so that
\begin{equation}\label{eq:ra-bv-buffer}
 B_L,\ B_{\mathrm{tail}},\ B_0\ge A+72.
\end{equation}
\begin{lemma}[The three low discrepancies]\label{lem:ra-n4-bv}
Under \eqref{eq:ra-bv-buffer},
\begin{equation}\label{eq:ra-three-bv-errors}
 |D_{I_L}(b_{\mathrm{all}})|+
 |D_{I_H}(b_{\mathrm{tail}})|+
 |D_{I_H}(c_{\mathrm{lo}})|\ll_A NL^{-A-10}.
\end{equation}
\end{lemma}
\begin{proof}
The arrays \(b_{\mathrm{all}}\), \(b_{\mathrm{tail}}\) and \(c_{\mathrm{lo}}\)
are supported on \(q<R_L\), \(q<R_{\mathrm{tail}}\) and \(q\le Q_0\),
respectively. Each is one-bounded. Since
\(B_L,B_{\mathrm{tail}},B_0\ge A+72>S+60\), every such \(q\) has
\(q\le N^{1/2}L^{-(S+60)}\).
This is the level of \eqref{eq:ra-prime-bv} at
\(X=N\). Every supported \(q\) is coprime to \(N\), so \(N\bmod q\) is a
reduced class. Since \(\pi(y;q,N)=\operatorname{Li}(y)=0\) for \(y<2\),
\[
 \Pi_{I_L}(q)-\frac{\operatorname{Li}(I_L)}{\varphi(q)}
   =\Delta_\pi(N/4;q,N).
\]
Write \(y^-\) for the left limit at \(y\). Then
\[
 \Pi_{I_H}(q)-\frac{\operatorname{Li}(I_H)}{\varphi(q)}
   =\Delta_\pi\bigl(((1-\varepsilon)N)^-;q,N\bigr)
    -\Delta_\pi(N/4;q,N).
\]
A left limit is a limit of values at points of \([0,N]\).
Hence each value on the right is at most
\(\sup_{0\le y\le N}|\Delta_\pi(y;q,N)|\).
Therefore the left side of
\eqref{eq:ra-three-bv-errors} is at most five times the left side of
\eqref{eq:ra-prime-bv}, with \(X=N\) and \(S=A+10\). This is the maximal
prime statement of Fact~\ref{fact:prime-estimates-maximal}.
\end{proof}


\subsection{The high coefficient}

Put
\begin{equation}\label{eq:ra-rough-arrays}
\begin{gathered}
 U=V=N^{1/3},\qquad z_0=N^{1/(\log\log N)^3},\\
 \mathcal R(n)=\mathbf1_{P^-(n)>z_0},\qquad \mathcal R(1)=1,\qquad
 \mathcal A(a)=\mu(a)\mathcal R(a),\\
 \mathcal C(t)=\mathcal R(t)
       \sum_{\substack{ab=t\\a\le U,\ b\le V}}\mu(a)\Lambda(b),\qquad
 \mathcal B(s)=\mathcal R(s)\sum_{\substack{b\mid s\\b>V}}\Lambda(b),\\
 \mathcal K_q(e)=\mathbf1_{e\equiv N\pmod q}
                   -\frac{\mathbf1_{(e,q)=1}}{\varphi(q)}.
\end{gathered}
\end{equation}
The von Mangoldt identity gives three sums:
\begin{equation}\label{eq:ra-three-bilinear-sums}
\begin{aligned}
 T_1&=\sum_q\Gamma(q)
   \sum_{\substack{a\le U,\ c\ge1\\ac\in I_H}}
   \mathcal A(a)\mathcal R(c)\frac{\log c}{\log(ac)}\mathcal K_q(ac),\\
 T_2&=-\sum_q\Gamma(q)
   \sum_{\substack{t\le UV,\ c\ge1\\tc\in I_H}}
   \frac{\mathcal C(t)\mathcal R(c)}{\log(tc)}\mathcal K_q(tc),\\
 T_{\mathrm{II}}&=\sum_q\Gamma(q)
   \sum_{\substack{a>U,\ s>V\\as\in I_H}}
   \frac{\mathcal A(a)\mathcal B(s)}{\log(as)}\mathcal K_q(as).
\end{aligned}
\end{equation}
These come from
\begin{equation}\label{eq:ra-vaughan-identity}
 \Lambda=\mu_{\le U}*\log-\mu_{\le U}*\Lambda_{\le V}*1
             +\Lambda_{\le V}+\mu_{>U}*\Lambda_{>V}*1.
\end{equation}
Roughness is multiplicative, including shared prime factors.
The short \(\Lambda\)-term vanishes on \(I_H\).

The detector retains proper prime powers whose bases exceed \(z_0\).
For \(e\in I_H\), hence \(e>z_0\),
\begin{equation}\label{eq:ra-proper-power-detector}
 \frac{\mathcal R(e)\Lambda(e)}{\log e}
 =\mathbf1_{e\in\mathbb P}
  +\sum_{\substack{P\in\mathbb P,\ j\ge2\\e=P^j}}
                                      \frac{\mathbf1_{P>z_0}}j.
\end{equation}
Consequently
\begin{equation}\label{eq:ra-b10}
\begin{aligned}
 P_{I_H}(\Gamma)&=T_1+T_2+T_{\mathrm{II}}-E_{\mathrm{pp}},\\
 E_{\mathrm{pp}}&=\sum_q\Gamma(q)
   \sum_{\substack{P\in\mathbb P,\ j\ge2\\P^j\in I_H}}
           \frac{\mathbf1_{P>z_0}}j\mathcal K_q(P^j).
\end{aligned}
\end{equation}
Let \(e\in I_H\).
A supported \(q\) with \(e\equiv N\pmod q\) divides
\(N-e\), and \(|\Gamma(q)|\le1\) by unique marking.
The comparator part is at most \(\sum_{q\le Q}1/\varphi(q)\). Hence \eqref{eq:ra-totient-mass}
gives
\begin{equation}\label{eq:ra-proper-power-bound}
 \sum_q|\Gamma(q)||\mathcal K_q(e)|\le d(N-e)+O(L),
 \qquad E_{\mathrm{pp}}\ll_\nu N^{1/2+\nu}L^2.
\end{equation}
Here \(N-e\ge\varepsilon N\), and \(0<\nu<1/2\) is fixed. The second bound
follows from the first, \(d(N-e)\ll_\nu N^{\nu}\) and
\(\sum_{j\ge2}\#\{P\in\mathbb P:P^j\le N\}\ll N^{1/2}\).
The high-coefficient estimate uses the \(P>z_0\) mask of
\(E_{\mathrm{pp}}\) in \eqref{eq:ra-b10}. That mask excludes bases
\(P\le z_0\),
so the displayed decomposition needs no separate add-back.

Proposition~\ref{hyp:ra-applications}(3) proves the high-coefficient estimate
\begin{equation}\label{eq:ra-ra1}
 P_{I_H}(\Gamma)\ll_{A,B,B_0}NL^{-A}
 \qquad(A>0\ \hbox{fixed}).
\end{equation}
Its modulus coefficient is exactly \eqref{eq:ra-factories}.
The estimate permits no joint arithmetic mask.
Table~\ref{tab:ra-rows} gives its decomposition.
Each power boundary belongs to the interval on its right.

\begin{table}[ht]
\caption{The high-coefficient decomposition}\label{tab:ra-rows}
\small
\begin{tabular}{@{}lll@{}}
\hline
Sum and range & Tested array & Required estimate\\
\hline
\(T_1\), all \(a\le U\) & \(\mathcal R(c)\) & Rough-factor estimate\\
\(T_2\), \(c<N^{.46}\) & \(\mathcal R(c)\) & Two-factor convolution\\
\(T_2\), \(N^{.46}\le c<N^{.50}\) & \(\mathcal C(t)/(2L)\) & Central estimate\\
\(T_2\), \(N^{.50}\le c<N^{.54}\) & \(\mathcal R(c)\) & Central estimate\\
\(T_2\), \(c\ge N^{.54}\) & \(\mathcal R(c)\) & Rough-factor estimate\\
\(T_{\mathrm{II}}\), \(a<N^{.46}\) & \(\mathcal A(a)\) & Two-factor convolution\\
\(T_{\mathrm{II}}\), \(N^{.46}\le a<N^{.50}\) & \(\mathcal B(s)/(2L)\) & Central estimate\\
\(T_{\mathrm{II}}\), \(N^{.50}\le a<N^{.54}\) & \(\mathcal A(a)\) & Central estimate\\
\(T_{\mathrm{II}}\), \(a\ge N^{.54}\) & \(\mathcal B(s)/(2L)\) & Two-factor convolution\\
\hline
\end{tabular}
\end{table}


\subsection{The three distribution estimates}

\begin{lemma}[Rough-factor distribution under parent-scale and strict-roughness conditions]\label{hyp:ra-rough}
Let \(L=\log N\), \(z_0=N^{1/(\log\log N)^3}\), and
\(Q=N^{5227/10000}L^{-B}\), with fixed \(B\). Assume the two explicit
conditions used by the direct rows:
\begin{enumerate}
 \item[(W1)] The original parent satisfies
 \(cN^{.54}\le H\le N\), for a fixed \(c>0\). In the \(T_1\) row
 \(H=N/a\le N\); in the right \(T_2\) row \(H=N/t\le N\).
 \item[(W2)] Use the strict roughness convention
 \(\mathcal R(n)=\mathbf1_{P^-(n)>z_0}\), as in
 \eqref{eq:ra-rough-arrays}. It is not replaced by
 \(P^-(n)\ge z_0\).
\end{enumerate}
For every fixed \(T_{\mathrm{save}}>0\), uniformly in real intervals \(I\subset(0,2H]\)
(so \(|I|\le2H\); open, closed or one-point),
supported \(q\le Q\) with \((q,N)=1\), and reduced residues \((b,q)=1\),
\begin{equation}\label{eq:ra-theorem-r}
 \sum_{n\in I}\mathcal R(n)
 \left\{\mathbf1_{n\equiv b\pmod q}
          -\frac{\mathbf1_{(n,q)=1}}{\varphi(q)}\right\}
 \ll_{T_{\mathrm{save}}}\frac{H}{qL^{T_{\mathrm{save}}}}.
\end{equation}
The interval may be arbitrarily short.
The saving is measured at its original parent. The estimate is an integer-sieve statement and uses no prime
distribution input.
\end{lemma}
\begin{proof}
Put \(\ell=\log\log N\) and
\(\mathcal P_q=\{p\le z_0:p\nmid q\}\). The Euler product at
\(s=1+1/\log z_0\) gives
\(h_q=\sum_{p\in\mathcal P_q}p^{-1}\le3\ell\) for sufficiently large
\(N\). Choose the fixed constant \(K\) after \(T_{\mathrm{save}}\), with
\begin{equation}\label{eq:ra-bonferroni-constant}
 K>3e,\qquad K\log\!\left(\frac K{3e}\right)>T_{\mathrm{save}}+4.
\end{equation}
Let \(k=2\lceil K\ell/2\rceil\). The two adjacent Bonferroni truncations
bracket the strict exclusion of all primes in \(\mathcal P_q\).
Their normalized gap satisfies
\begin{equation}\label{eq:ra-theorem-r-lattice}
 k=2\lceil K\log\log N/2\rceil,\qquad
 E_k\le\frac{h_q^k}{k!}
       \le\left(\frac{3e}{K}\right)^{K\ell}
       \le L^{-T_{\mathrm{save}}-4},\qquad
 (k+1)z_0^k=N^{O(K/\ell^2)}.
\end{equation}
Every divisor in the truncation is at most \(z_0^k\).
By CRT, each divisor
contributes its expected interval mass with error \(O(1)\).
This is uniform for open, closed, or singleton intervals.
The same expansion applies to the comparator after using
\[
 \mathbf1_{(n,q)=1}=\sum_{\substack{d\mid\operatorname{rad}(q)\\d\mid n}}\mu(d).
\]
Thus the main terms cancel. The accumulated lattice error is
\(N^{O(K/\ell^2)}q/\varphi(q)\), and
\(q/\varphi(q)\le\omega(q)+1\ll L\). Under (W1),
\(H/q\gg N^{.0173}L^B\). The Bonferroni remainder is at
most \((|I|/q)E_kL^2\le2HL^{-T_{\mathrm{save}}-2}/q\) because \(|I|\le2H\). Hence both are
absorbed by \(H/(qL^{T_{\mathrm{save}}})\) once the fixed truncation precision is raised. This proves \eqref{eq:ra-theorem-r}.

The strict convention in (W2) has no threshold atom.
An application might instead use \(P^-(n)\ge z_0\).
Then the only possible difference is the atom
\(z_0=P\) when \(z_0\) is prime. That weak convention requires its separate
CRT boundary correction. It is not used in \eqref{eq:ra-theorem-r}.
\end{proof}
The positive gap \(.54-v_{\mathrm{RA}}=173/10000\) absorbs the lattice error
in \eqref{eq:ra-theorem-r-lattice}.
For \(T_1\), the free scale is at least \(N^{2/3}/4\).
For the right \(T_2\)-part, it exceeds a fixed multiple of \(N^{.54}\).
Partial summation then gives
\begin{equation}\label{eq:ra-direct-row-total}
 \sum_{m\le M_{\mathrm{out}}}\sum_{q\le Q}\frac{N}{mqL^{T_{\mathrm{save}}}}
 \ll NL^{2-T_{\mathrm{save}}}.
\end{equation}
Both kernel terms remain in this estimate.
Here \(M_{\mathrm{out}}\) is the outer endpoint of the current row.
It is \(U\) in the \(T_1\) row and the
corresponding low-row endpoint in the \(T_2\) decomposition.


The other tested arrays need uniformity under coprimality
with an auxiliary integer. Write \(\tau(h)=\sum_{d\mid h}1\) and
\(\omega(h)=\#\{p\in\mathbb P:p\mid h\}\).
For any sequence \(f(n)\) on its original parent \((H,2H]\), define
\begin{equation}\label{eq:ra-sw-difference}
\begin{aligned}
 \Delta_f(I;k,b;h)={}&
 \sum_{\substack{n\in I,\ n\equiv b\pmod k\\(n,h)=1}}f(n)\\
 &-\frac1{\varphi(k)}
     \sum_{\substack{n\in I\\(n,kh)=1}}f(n) .
\end{aligned}
\end{equation}
\begin{lemma}[Prime-variable transfer]\label{lem:ra-prime-transfer}
Let \(N^{1/3}/8\le H\le 2N\), \(L=\log N\), and
\(z_0=N^{1/(\log\log N)^3}\). Let \(y_0\ge z_0\) and
\[
 f(n)=\sum_{\substack{mp=n\\p\in J_m}}\alpha(m)w(p),
\]
where \(p\) is prime and \(|\alpha(m)|\le1\).
For each \(m\) the
allowed primes lie in a real interval \(J_m\subset[y_0,\infty)\).
Assume \(\|w\|_\infty+\operatorname{Var}(w)\le W\) on \([y_0,2H]\)
and \(|f(n)|\le1\) on \((H,2H]\). For every fixed \(P>0\), uniformly
in intervals \(I\subset(H,2H]\), \(k,h\ge1\), and \((b,k)=1\),
prime Siegel--Walfisz, Paper~I (2.21), implies
\begin{equation}\label{eq:ra-prime-transfer}
 |\Delta_f(I;k,b;h)|
 \ll_P WH\left(L^{-P}+\frac{\omega(h)L}{y_0}\right)
 \ll_P WH\tau(h)L^{-P}.
\end{equation}
\end{lemma}
\begin{proof}
Put \(K=2P+4\). If \(k>L^K\), the two terms in the discrepancy are
bounded by a constant times \(\|f\|_\infty(H/k+1+H/\varphi(k))\).
For each represented \(n\), a prime divisor \(p\ge y_0\) determines at
most one pair \((m,p)\), since \(m=n/p\). Hence
\[
 |f(n)|\le W\frac{\log(2H)}{\log y_0}
 \ll W(\log L)^3,
\]
using \(y_0\ge z_0\), \(H\le2N\), and
\(\log z_0=L/(\log L)^3\). Since \(\varphi(k)\gg\sqrt{k}\) and
\(H\ge N^{1/3}/8\), the resulting discrepancy is
\(\ll_P WHL^{-P}\). The two spare powers in \(K/2=P+2\) absorb
\((\log L)^3\). The additive \(1\) is absorbed because
\(H\ge N^{1/3}/8\gg L^P\).

Now let \(k\le L^K\). Terms with \((m,kh)>1\) vanish from both sides.
Otherwise the prime progression is \(p\equiv bm^{-1}\pmod{k}\) on
\(J_m\cap I/m\).
Primes dividing \(h\) are excluded on the progression side and in the
comparator.
For each \(m\) the two exclusions differ by at most
\(2W\omega(h)\).
There are at most \(2H/y_0\) possible
values of \(m\), so these terms contribute
\(\ll WH\omega(h)/y_0\).

For the remaining primes define
\[
 E(t)=\pi(t;k,c)-\frac1{\varphi(k)}
       \sum_{\substack{a\bmod k\\(a,k)=1}}\pi(t;k,a).
\]
The reduced-class sum equals \(\pi(t)-\#\{p\le t:p\mid k\}\).
Thus this is the centered prime discrepancy without the primes dividing \(k\).
Apply (2.21) to the class \(c\) and average its uniform error over
all reduced classes.
This gives \(E(t)\ll_{A'}t(\log t)^{-A'}\), with no factor
of \(k\). Since \(t\ge y_0\ge z_0\),
\(\log t\ge L/(\log L)^3\). Hence, for all sufficiently large \(N\),
\(k\le(\log t)^{2K}\).
Therefore (2.21), with modulus exponent \(2K\)
and fixed saving \(A'\), applies throughout each interval of
integration. Stieltjes partial summation bounds each \(m\)-contribution
by \(O(W(H/m)(\log y_0)^{-A'})\). Sum over
\(m\le2H/y_0\) and take \(A'=2(P+K+2)\).
This gives \(O_P(WHL^{-P})\). The estimate is measured at the original parent
\(H\), including for short intervals. Finally,
\(L/y_0\ll_P L^{-P}\) because \(y_0\ge z_0\), and
\(\omega(h)\le\tau(h)\). This proves \eqref{eq:ra-prime-transfer}.
\end{proof}

\begin{lemma}[Four-family auxiliary-coprime Siegel--Walfisz]
\label{lem:ra-sw}
By prime Siegel--Walfisz, Paper~I (2.21), for each of
\(\mathcal R,\mathcal A,\mathcal C/(2L),\mathcal B/(2L)\), the bound
\begin{equation}\label{eq:ra-full-sw}
 |\Delta_f(I;k,b;h)|\ll_P H\tau(h)L^{-P},
 \qquad (b,k)=1,
\end{equation}
holds uniformly for every \(I\subset(H,2H]\) and \(k,h\ge1\).
The parents are the
central ranges in Table~\ref{tab:ra-rows} and the convolution parents
\(N^{1/3}/8\le H\le2N^{.46+O(1/L)}\).
The proof uses the lower
bound only to absorb the large-modulus lattice term and to exclude \(n=1\).
It uses \(H\le2N^{.54}\) for the remaining weight bounds. The divisor
exponent is the single fixed value \(C_{\mathrm{SW}}=1\), independent of
\(P\). The implied constant may depend on \(P\) and the fixed family
parameters, but not on \(N,H,I,k,h,b\). The error is measured at \(H\),
even when \(I\) is short.
\end{lemma}
\begin{proof}
For sufficiently large \(N\), every parent has
\(H\ge N^{1/3}/8>1\), so \(n=1\) is excluded. The rough and signed-rough
arrays have exact prime-variable decompositions.
For \(\mathcal R(n)\), take \(\alpha(m)=\mathcal R(m)\),
\(J_m=[P^+(m),\infty)\cap(z_0,\infty)\), and \(w(p)=1\).
Use \(P^+(1)=1\). The largest prime factor gives the unique term. For
\(\mathcal A(n)=\mu(n)\mathcal R(n)\),
take \(\alpha(m)=-\mu(m)\mathcal R(m)\),
\(J_m=(\max(P^+(m),z_0),\infty)\), and \(w(p)=1\). The strict
inequality excludes repeated largest primes, where \(\mu(n)=0\).

For \(\mathcal B(s)/(2L)\), the \(j=1\) terms in
\(b=p^j\mid s\), \(b>V\), use \(\alpha(m)=\mathcal R(m)\),
\(J_m=(V,\infty)\), and \(w(p)=\log p/(2L)\). For
\(\mathcal C(t)/(2L)\), the \(j=1\) terms use
\(\alpha(a)=\mu(a)\mathcal R(a)\mathbf1_{a\le U}\),
\(J_a=(z_0,V]\), and the same weight. In each case
\(\|w\|_\infty+\operatorname{Var}(w)\le1\).
The four tested arrays
are one-bounded on their parent ranges. For the last two arrays this follows
from \(\sum_{d\mid n}\Lambda(d)=\log n\) and
\(\log(2H)<2L\). On the central ranges \(H\le2N^{.54}\).
There the weight bound follows from \(y_0=V=N^{1/3}\) for \(\mathcal B\) and \(y_0=z_0\)
for \(\mathcal C\), since \(\log(2H)\le .54L+O(1)\).

For \(\mathcal B\) and \(\mathcal C\), write the remaining \(j\ge2\)
terms as \(f_{\ge2}\). They need not be nonnegative.
For instance, the
\(\mu(a)\) coefficient in \(\mathcal C\) may change sign.
Take absolute values before summing:
\[
 \sum_{n\le2H}|f_{\ge2}(n)|
 \ll \frac{H}{L}\sum_{p>z_0}\sum_{j\ge2}\frac{\log p}{p^j}
 \ll \frac{H\log z_0}{Lz_0}
 \ll \frac{H}{z_0}.
\]
For any sequence \(g\),
\(|\Delta_g|\le2\sum_{n\le2H}|g(n)|\).
Hence this remainder is
\(O_P(HL^{-P})\) for every fixed \(P\). Applying
Lemma~\ref{lem:ra-prime-transfer} to each prime term proves
\eqref{eq:ra-full-sw} with the same exponent \(C_{\mathrm{SW}}=1\).
\end{proof}
The mean-square estimate below has five uses.
Section~\ref{sec:k2} uses it for the zero mode of the exterior sum.
Section~\ref{sec:retained-side} uses it for the Class II estimate of
Paper~II inside the retained side.
Section~\ref{sec:l3} uses it for the row-three dispersion estimate of the
switched sum. These three apply Paper~I, Theorem~3.2.
In the large-norm case of Paper~I, (3.48), part (ii) below replaces
\cite[Theorem~0(a)]{BFI86}. The other two uses are the central rectangular
estimate of Lemma~\ref{hyp:ra-central} and the small-modulus norm-relative
step of Lemma~\ref{hyp:ra-l52}. Part (i) derives the conclusion of that
published theorem (Paper~I, Theorem~2.4).
Its inputs are the norm-relative condition
\eqref{eq:ra-mean-square-a2} and
an elementary large sieve, so the argument is self-contained
at this step. We do not correct \cite{BFI86}.
Paper~I's use of it is unaffected. The norm-relative premise is not deduced from the
absolute parent-scale estimate \eqref{eq:ra-full-sw} alone.
Let \(\xi\) be complex and supported in \([1,Z]\), \(Z\ge3\). Fix
\(B_{\mathrm{div}}\ge0\). Consider the following condition.
For every
\(P>0\), uniformly in integers \(h,k\ge1\) and reduced classes
\(b\bmod k\),
\begin{equation}\label{eq:ra-mean-square-a2}
 \left|\sum_{\substack{n\equiv b\ (k)\\(n,h)=1}}\xi_n
 -\frac1{\varphi(k)}\sum_{(n,kh)=1}\xi_n\right|
 \ll_P\|\xi\|_2Z^{1/2}\tau(h)^{B_{\mathrm{div}}}(\log(2Z))^{-P},
\end{equation}
The constants are common to the family of pairs \((Z,\xi)\) under
consideration.
Write \(\Delta_\xi(k,b)\) for the centered discrepancy at \(h=1\). The norm conclusion is
\begin{equation}\label{eq:ra-transferred-mean-square}
 \sum_{k\le Z(\log Z)^{-B_1(A)}}\sum_{b\bmod k}^{*}
       |\Delta_\xi(k,b)|^2
 \ll_A\|\xi\|_2^2 Z(\log Z)^{-A}.
\end{equation}
In the selected-factor form of Paper~I, (3.48),
\(\xi_n=\alpha_n1_{(n,r)=1}\), \(Z=\lceil2T_{\mathrm{int}}\rceil\), and
\begin{equation}\label{eq:ra-mean-square-transfer}
 \sum_{q\le3Q}\sum_{b\bmod q}^{*}|\Delta_\xi(q,b)|^2
 \ll N_{\mathcal I}T_{\mathrm{int}} L^{-A_3},\qquad
 3Q\le Z(\log Z)^{-B_F},\quad B_F=A_3+15,
\end{equation}
where \(N_{\mathcal I}\) is the interval mass in Paper~I, (3.8).
\begin{lemma}[Transferred mean-square estimate]
\label{hyp:ra-ms-transfer}
\begin{enumerate}
 \item[(i)] If \eqref{eq:ra-mean-square-a2} holds, then
 \eqref{eq:ra-transferred-mean-square} holds for every fixed \(A>0\).
 Here \(B_1(A)\) and the implied constant depend only on \(A\),
 \(B_{\mathrm{div}}\) and the common constants.
 \item[(ii)] Assume the hypotheses (3.1), (3.2), (3.6), (3.8) with its
 length condition, (3.11) and (3.25)--(3.27) of Paper~I. Then the estimate
 \eqref{eq:ra-mean-square-transfer} holds uniformly over the parameter
 boxes and selected factors of Paper~I, Theorem~3.2. Hence that theorem
 holds with this input in place of \cite[Theorem~0(a)]{BFI86} at (3.48).
\end{enumerate}
\end{lemma}
\begin{proof}
(i) Expand \(\Delta_\xi(k,b)\) in Dirichlet characters modulo \(k\) and
pass to primitive characters. For conductors up to a fixed power of
\(\log Z\), apply \eqref{eq:ra-mean-square-a2}.
Take the auxiliary divisor
\(h\) to be the imprimitive part of the modulus. For larger conductors
use the multiplicative large sieve inequality.
Its elementary proof uses
finite Fourier identities and the spacing of reduced fractions.
Summing over the moduli gives (i).
(ii) If \(\|\xi\|_2^2\le N_{\mathcal I}L^{-A_3-1}\), the elementary bound
in Paper~I, (3.48), applies. Otherwise (3.25), with auxiliary divisor
\(dr\), and (3.26) give \eqref{eq:ra-mean-square-a2} for
\(\xi_n=\alpha_n1_{(n,r)=1}\) with divisor exponent one. Also (3.6) gives
the modulus range in \eqref{eq:ra-mean-square-transfer}. Part (i) then
applies. The complete argument is in Appendix~
\hyperref[subsec:supp-lem-65]{\ref*{subsec:supp-lem-65}}, with its
row index in Appendix~
\hyperref[subsec:app-ra-rows]{\ref*{subsec:app-ra-rows}}. We assert no implication from
\eqref{eq:ra-full-sw} alone to the norm bound.
\end{proof}

For the three convolution rows, let
\begin{equation}\label{eq:ra-l52-scales}
 D_1=HN^{-\sigma_{\mathrm p}},\qquad
 D_2=QN^{w_{\mathrm{RA}}+\sigma_{\mathrm p}}/H,\qquad
 \sigma_{\mathrm p}=10^{-6},\quad w_{\mathrm{RA}}=9/1000,
\end{equation}
where \(N^{1/3}/8\le H\le2N^{.46+O(1/L)}\).
A supported modulus has
the unique form \(q=pd\) of \eqref{eq:ra-coefficient}.
Here \(p\) is the marked prime and \(d\) the squarefree cofactor.
\begin{equation}\label{eq:ra-l52-caps}
 r<D_1,\qquad s<D_2,\qquad rs=q<Q.
\end{equation}

\begin{lemma}[Deterministic allocation]\label{lem:ra-selector}
For each supported \(q=pd\), the following rule defines a unique coprime
split \(q=rs\) satisfying \eqref{eq:ra-l52-caps}. Put \(p\) in \(r\).
Order the cofactor primes exceeding \(N^{w_{\mathrm{RA}}}\) as
\(b_1>\cdots>b_m\).
Here \(m\le(v_{\mathrm{RA}}-t_0)/w_{\mathrm{RA}}=37553/1260<30\). Process them in this order.
Put \(b_i\) in \(r\) exactly when the current \(r b_i<D_1\), and put it
in \(s\) otherwise. Let \(r_{\rm big}\) and \(s_{\rm big}\) be the
resulting factors.
Order the remaining cofactor primes as
\(l'_1>\cdots>l'_k\). Set
\(r_0=r_{\rm big}\prod_{j=1}^k l'_j\). If \(r_0<D_1\), take
\(r=r_0\), \(s=s_{\rm big}\). Otherwise let \(j_*\) be the least index
for which \(r_0/(l'_1\cdots l'_{j_*})<D_1\), and take
\[
 r=r_0/(l'_1\cdots l'_{j_*}),\qquad
 s=s_{\rm big}l'_1\cdots l'_{j_*}.
\]
Call the resulting pair \(\operatorname{Sel}_H(q)=(r,s)\).
\end{lemma}
\begin{proof}
The marked prime fits, since
\(p<N^{1/3-\eta_{\rm m}}<D_1\), even at the fixed lower endpoint
\(H\ge N^{1/3}/8\). The accepted big-prime partial products stay below
\(D_1\).
Thus \(r_{\rm big}<D_1\), and the least index \(j_*\) exists
when needed. Write all cofactor primes in decreasing order as
\(l_1>l_2>\cdots\). The odd cube gates in \eqref{eq:ra-coefficient}
imply
\[
 p l_1\cdots l_{i-1}l_i^2<Q\qquad\text{for every }i:
\]
For odd \(i\), drop one factor from the cube gate.
For even \(i\), use
the preceding odd gate and \(l_i\le l_{i-1}\).

Suppose \(r_0<D_1\), and let \(b_i\) be the last big prime rejected
into \(s\). The rejection test gives \(r^{(i)}b_i\ge D_1\). The
two-gate inequality then gives
\[
 s=s_{\rm big}
 =\frac{p b_1\cdots b_i}{r^{(i)}}
 \le\frac{p b_1\cdots b_{i-1}b_i^2}{D_1}<\frac Q{D_1}<D_2.
\]
If no big
prime is rejected, \(s=1\). In the prefix case, minimality gives
\(r l'_{j_*}\ge D_1\), hence \(r\ge D_1N^{-w_{\mathrm{RA}}}\) and
\[
 s=q/r<QN^{w_{\mathrm{RA}}}/D_1=D_2.
\]
The factors partition the squarefree prime factors of \(q\).
Thus
\((r,s)=1\), \(rs=q<Q\), and \(r<D_1\).
\end{proof}

Define the exact joint coefficient on each dyadic Cartesian modulus shell by
\begin{equation}\label{eq:ra-global-tensor}
\begin{gathered}
 F_H(r,s)=
 \begin{cases}
 \Gamma(rs),&rs\text{ is supported and }\operatorname{Sel}_H(rs)=(r,s),\\
 0,&\text{otherwise},
 \end{cases}\\
 \sum_q\Gamma(q)\Phi(q)=\sum_{r,s}F_H(r,s)\Phi(rs),\\
 F_H(r,s)=\sum_\nu c_\nu f_\nu(r)g_\nu(s),\qquad
 |f_\nu|,|g_\nu|\le1,\qquad
 \sum_\nu|c_\nu|\ll L^{C_{\mathrm{tensor}}},\quad C_{\mathrm{tensor}}=2^{30}+91.
\end{gathered}
\end{equation}
Here \(\operatorname{Sel}_H\) is the rule in Lemma~\ref{lem:ra-selector}.
The expansion is exact on the full Cartesian parent, including zero values
off the selected pairs.

\begin{lemma}[Global tensor for the selector]\label{lem:ra-tensor}
The coefficient \(F_H\) has a rank-one expansion satisfying
\eqref{eq:ra-global-tensor}, with the sharper bound
\begin{equation}\label{eq:ra-tensor-norm}
 \sum_\nu|c_\nu|
 \le4(2^{30}-1)(3+L)^{78}\ll L^{79}.
\end{equation}
\end{lemma}
\begin{proof}
Fix a pair of dyadic shells. A branch fixes the merge order
\(\pi\in\{R,S\}^m\) of the at most 29 big cofactor primes. It also fixes the final
large-prime/prefix case.
In the prefix case it fixes the parity of the number
of factors assigned to \(s\). There are at most
\(4\sum_{m=0}^{29}2^m=4(2^{30}-1)\) branches. Within a branch, unary
conditions impose squarefreeness, coprimality with \(2N\) and the unique
marked prime \(p\) in \(r\). They also impose the bounds \(P^+(r/p)<z_s\) and
\(P^+(s)<z_s\), assigned big-prime counts, \(r<D_1\), \(s<D_2\), and
\(r_{\rm big}<D_1\). The remaining unary data are the case indicator, the parity of the small-prime
count, and the sign \(\Gamma(rs)=\mu(rs/p)=-\mu(r)\mu(s)\).

Every cross-condition is a comparison \(\phi(r)<\psi(s)\) or its weak
counterpart. Let \(K_0\) be the number of distinct values of \(\phi\) on
the \(r\)-shell. For ranks \(1\le a\le K_0\), \(0\le b\le K_0\), put
\(M=2K_0+1\) and
\[
 c_t=M^{-1}\sum_{u=0}^{K_0}e(-tu/M).
\]
Finite Fourier orthogonality gives the exact separator
\begin{equation}\label{eq:ra-rank-separator}
 \mathbf1_{a\le b}=\sum_{t\bmod M}c_t e(tb/M)e(-ta/M),
 \qquad \sum_{t\bmod M}|c_t|\le3+\log M.
\end{equation}
Rank \(\phi(r)\) among its shell values, and count the values below
\(\psi(s)\).
Each shell has at most \(Q\) values, so
\(\log M\ll L\). The cross-comparisons are exactly:
\begin{enumerate}
 \item the adjacent strict comparisons in the fixed merge order of the
 big primes (at most 28);
 \item for each \(S\)-letter \(b_i\), the greedy rejection test
 \(p(\prod_{b_j>b_i,\ b_j\in r}b_j)b_i\ge D_1\) (at most 29);
 \item the odd cube gates
 \(p l_1\cdots l_{i-1}l_i^3<Q\) at odd indices among the big primes
 (at most 15);
 \item the remaining odd cube gates on the small-prime part, grouped by
 the fixed parity into the two maxima of the \(r\)- and \(s\)-side
 partial products (two comparisons);
 \item in the prefix case, the prefix-order test
 \(\min P>\max R_{\rm small}\) and the least-prefix test
 \(r\min P\ge D_1\) (two comparisons);
 \item the support range \(Q_0<rs<Q\) (two comparisons).
\end{enumerate}
There are at most 78 comparisons in total. Once a merge branch is fixed,
each displayed test compares an \(r\)-only quantity with an
\(s\)-only quantity. The small-prime maxima combine all gates of the same
parity and side. The branch conditions are exactly the greedy tests, the
cube gates and the least-prefix condition.
Hence their union is precisely
\(F_H\). Multiplying the separators in
\eqref{eq:ra-rank-separator} gives \eqref{eq:ra-tensor-norm}. The unary
factors remain one-bounded. The bound in \eqref{eq:ra-tensor-norm} is per
dyadic shell pair. The resulting \(O(L^2)\) shell pairs are absorbed by the
chosen logarithmic precision.
\end{proof}

For the convolution rows, split \(s=dt\) with
\(P^+(d)\le z< P^-(t)\), where
\(z=N^{1/(\log\log N)^3}\) and \(y=N^{1/\log\log N}\). The aggregate
with \(d>y\) is bounded by Rankin's method. The mean-square modulus in the
proof of Lemma~\ref{hyp:ra-l52} is \(k=rd\), not the full \(q\).
Its scale is
\begin{equation}\label{eq:ra-l52-small-modulus}
 Q'=D_1y=HN^{-\sigma_{\mathrm p}+1/\log\log N},\qquad
 k\le5Q',\qquad Z=\lceil2H\rceil.
\end{equation}
The inner factor \(t\) is \(z\)-rough and is removed in \(L^2\), with cost
\(H^2L^C/z\). No bound for \(\tau(r)\) and no roughness condition on \(r\)
is used. The required range is
\begin{equation}\label{eq:ra-l52-range}
 5Q'\le Z(\log Z)^{-B_F},\qquad B_F=P_{\mathrm{MS}}+15,
\end{equation}
where \(P_{\mathrm{MS}}\) is the fixed source-level precision in the
supplementary Lemma~6.5(i). This holds since
\(Q'/H=N^{-\sigma_{\mathrm p}+o(1)}\). The full-modulus
range \(3Q\le Z(\log Z)^{-B_F}\) is not asserted for these parents.

\begin{lemma}[Two-factor convolution estimate]\label{hyp:ra-l52}
Let \(N^{1/3}/8\le H\le2N^{.46+O(1/L)}\) and \(HM\asymp N\). Let \(a_n\),
\(n\sim H\), be one-bounded, supported on \(P^-(n)>z\), and satisfy
\eqref{eq:ra-full-sw} at the original parent \(H\), where
\(z=N^{1/(\log\log N)^3}\). Let \(b_m\), \(m\sim M\), be one-bounded.
Let \(U_1,V_1\) be dyadic modulus scales, allowing the singleton scale
\(1/2\) for the integer \(1\), such that
\[
 U_1\le D_1,\qquad V_1\le D_2,\qquad U_1V_1\le N^{v_{\mathrm{RA}}}.
\]
Let \(f(r)\), \(r\sim U_1\), and \(g(s)\), \(s\sim V_1\), be arbitrary
one-bounded arrays.
Keep the unary restrictions \(r<D_1\), \(s<D_2\),
\((r,N)=(s,N)=1\). Neither a pointwise condition \(rs<Q\) nor a joint condition
\((r,s)=1\) is imposed on the pair. Then, for every fixed \(A>0\), uniformly in these
occupied-parent scales, \(f\), \(g\), unary interval cuts of \(a_n\) and
\(b_m\), and intervals \(I\subset[N/8,2N]\),
\begin{equation}\label{eq:ra-l52-input}
 \sum_{r,s}f(r)g(s)
 \sum_{\substack{n,m\\nm\in I}}a_n b_m\mathcal K_{rs}(nm)\ll_A NL^{-A}.
\end{equation}
\end{lemma}
\begin{proof}
We use the dispersion method of Paper~I, Theorem~3.2.
It runs on the
whole Cartesian shell \(r\sim U_1\), \(s\sim V_1\), with residue class
\(N\bmod rs\). Its occupied-parent scale conditions are
\(U_1\le HN^{-\sigma_{\mathrm p}}\),
\(V_1\le N^{v_{\mathrm{RA}}+w_{\mathrm{RA}}}/(HN^{-\sigma_{\mathrm p}})\), and
\(U_1V_1\le N^{v_{\mathrm{RA}}}\).
The second follows from the displayed
bound by \(D_2\) after dropping \(L^{-B}\). The tested sequence is
strictly \(z\)-rough. The unary supports carry \((rs,N)=1\) and the
individual caps. No pointwise product cap is needed for the rank-one
arrays.
After Cauchy's inequality in \(m\), the split \(s=dt\) above reduces the
modulus of the mean square to \(k=rd\) in the range
\eqref{eq:ra-l52-range}. The mean square is split at a fixed power of
\(L\) below its target. Below it the trivial bound suffices. Above it
\eqref{eq:ra-full-sw} gives \eqref{eq:ra-mean-square-a2} with a loss of
that power, so Lemma~\ref{hyp:ra-ms-transfer}(i) applies. The
prime Siegel--Walfisz input is Paper~I, (2.21). No joint condition \((r,s)=1\) is used, because \(r,s\) are not the
moduli of Paper~I, Theorem~2.8. Write \(s=dt\) as above and \(k=rd\). A
prime \(p\mid(r,s)\) with \(p\le z\) divides \(d\) and only enlarges \(k\).
The Chinese remainder step needs only that \(k,t_1,t_2\) are pairwise
coprime.
The smooth--rough split gives \((d,t_1t_2)=1\). A prime
\(p\mid(r,s)\) with \(p>z\) divides \(k\) and a rough \(t_i\). Such tuples
are the bad type \((k,t_1t_2)>1\). If \(p\mid t_2\), the second progression
\(kt_2\mid mn_2-N\) places \(n_2\) in one class modulo \(kp\). If
\(p\mid t_1\) only, exchange the two copies first; the absolute square
majorant is symmetric, and the same statement then applies. We therefore
remove these tuples before the Chinese remainder step.
The cost in the square is
\(\ll MHL^C[H/(zK)+1]\). After the factor \(KM\) of the
first Cauchy step this is
\(\ll NL^C(z^{-1/2}+(K/H)^{1/2})\).
That is \(O_A(NL^{-A})\), since
\(z=N^{1/(\log\log N)^3}\) and \(K/H\le N^{-\sigma_{\mathrm p}/2}\).
The zero-mode identity holds for all tuples, including those with gcd
failures. After the second Cauchy step, Paper~I, Theorem~2.8 bounds the nonzero
frequencies. It is applied with \(r_{\mathrm{int}}=tt'\), the product of two
\(z\)-rough inner variables, and \(s_{\mathrm{int}}=p_0\), the least prime not dividing
\(N\). Since \(p_0\ll\log N<z\), \((r_{\mathrm{int}},s_{\mathrm{int}})=1\), exactly as in Paper~I,
(3.82)--(3.83). This is the form of Paper~I, Theorem~2.8
(Sections~5.3 and~7.4) used here. Its statement allows arbitrary one-bounded
\(f,g\) with only \((rs,N)=1\). The product restriction
\(nm\in I\) is kept as in Paper~I, Theorem~3.2.
The finite product
partition used there also handles the endpoints.
Neither \eqref{eq:ra-mean-square-transfer} at the full modulus \(3Q\) nor
a norm-relative consequence of \eqref{eq:ra-full-sw} for small norms is
used. The complete argument is in Appendix~
\hyperref[subsec:supp-lem-68]{\ref*{subsec:supp-lem-68}}, with its
row index in Appendix~
\hyperref[subsec:app-ra-rows]{\ref*{subsec:app-ra-rows}}.
\end{proof}


The central rows use \(Y=N^y\), \(y=341/5000\).
Choose the largest divisor \(s\mid d\) with \(s\le Y\).
Since \(p<N^{1/3-\eta_{\mathrm m}}\) and \(q=pd>Q_0\),
\(d>Q_0/p\gg N^{1/6+\eta_{\mathrm m}}L^{-B_0}>Y\).
Hence a next divisor exists.
The cube gates control adjacent divisors and give
\begin{equation}\label{eq:ra-central-selector}
 \frac{Yq}{Q}<s\le Y,\qquad r=q/s<Q/Y.
\end{equation}
This allocation keeps a one-bounded joint modulus coefficient.
The required central estimate has
\begin{equation}\label{eq:ra-m53-range}
\begin{gathered}
 H_{\mathrm c}M_{\mathrm c}\asymp N,\qquad
 N^{.50}\ll H_{\mathrm c}\ll N^{.54},\\
 Q_1\ll N^{.4545},\qquad Q_2\ll N^{.0682}.
\end{gathered}
\end{equation}
\begin{lemma}[Central rectangular distribution]\label{hyp:ra-central}
Let \(x\) be even, \(H_{\mathrm c}M_{\mathrm c}\asymp x\), and let the
bounds in \eqref{eq:ra-m53-range} hold with \(N\) replaced by \(x\), with
fixed comparison constants. Let \(a_n,b_m\) be one-bounded on the complete
dyadic rectangle \(n\sim H_{\mathrm c},m\sim M_{\mathrm c}\). Let the
tested \(a_n\) satisfy \eqref{eq:ra-full-sw} at its original parent.
The divisor exponent is fixed, and the constants are uniform at every requested
precision. Write
\[
 \mathcal D(q,a)=\sum_{n\sim H_{\mathrm c},m\sim M_{\mathrm c}}a_nb_m
 \left(1_{nm\equiv a\ (q)}-\frac{1_{(nm,q)=1}}{\varphi(q)}\right).
\]
Then, for every fixed \(A>0\),
\begin{equation}\label{eq:ra-m53-input}
 \sum_{r\sim Q_1}\max_{(b,r)=1}
 \sum_{\substack{s\sim Q_2\\\mu(rs)^2=1,\ (rs,x)=1}}
 \max_{\substack{a\bmod rs,\ (a,rs)=1\\a\equiv b\ (r)}}
       |\mathcal D(rs,a)|\ll_A x(\log x)^{-A}.
\end{equation}
The constant is uniform over these families, rectangles, squarefree
modulus supports and unary interval restrictions of \(a_n\) and \(b_m\).
All errors are measured at the original parents. One outer residue
\(b\) is chosen before \(s\) varies.
In the application it is
\(N\bmod r\).
\end{lemma}
\begin{proof}
This is the rectangular dispersion theorem proved in Appendix~
\hyperref[subsec:supp-lem-69]{\ref*{subsec:supp-lem-69}} and indexed
in Appendix~\hyperref[subsec:app-ra-rows]{\ref*{subsec:app-ra-rows}}.
The restriction
\((rs,x)=1\) only removes nonnegative terms. Its proof opens the inner
modulus \(s\) by dispersion with the outer residue fixed. It uses the mean
square over all moduli of Lemma~\ref{hyp:ra-ms-transfer}(i), which gives
Paper~I, (2.17), under its hypothesis (2.16). The supplementary proof applies the mean square twice.
The second application is repeated over two
correlations. At each application a constant sequence
on the same parent is added. Its sign is chosen by the parallelogram
identity to make the tested sequence norm-large. The resulting sequence
and the constant sequence satisfy \eqref{eq:ra-mean-square-a2} uniformly.
Their squared norms are at most the original parent times a fixed
power of \(L\). Fix the requested logarithmic modulus buffer \(B_1\)
first.
Then choose the source-level parameter so that the moduli lie in the
range of \eqref{eq:ra-transferred-mean-square}. Apply part (i) to both
sequences,
subtract by linearity, and absorb that fixed power by increasing the
requested precision. This uses neither part (ii) nor a mean-square
conclusion for an arbitrary small-norm sequence from absolute
Siegel--Walfisz alone. The other inputs are Paper~I, (2.21) (prime
Siegel--Walfisz), (2.35) (Weil--Estermann)
and Lemma~3.4 (completed hyperbola count). It does not use Paper~I, Theorem~2.8. For the four tested
arrays of Table~\ref{tab:ra-rows}, Lemma~\ref{lem:ra-sw} supplies
\eqref{eq:ra-full-sw} with divisor exponent one.
\end{proof}
Lemma~\ref{hyp:ra-central} concerns complete dyadic rectangles and
one-bounded sequences. Lemma~\ref{lem:ra-product-region} below reduces the
four central rows, with their product cuts, to it.
The raw exponent checks are
\begin{equation}\label{eq:ra-central-margins}
\begin{aligned}
 7v_{\mathrm{RA}}+5y&=3.9999<4,\\
 3v_{\mathrm{RA}}-y-1&=.4999<.50,\\
 1-v_{\mathrm{RA}}+y&=.5455>.54.
\end{aligned}
\end{equation}
The three margins verify the numerical range only. They do not establish
the rectangular-support or family hypotheses for a selected product region.


Product cuts require an additional finite reduction. For \(e\in I_H\), put
\begin{equation}\label{eq:ra-restoration}
\begin{aligned}
 w(e)&=\frac{2L}{\log e},\\
 \mathcal A(a)\mathcal R(c)\frac{\log c}{\log(ac)}
 &=\left(\mathcal A(a)\mathcal R(c)\frac{\log c}{2L}\right)w(ac),\\
 -\frac{\mathcal C(t)\mathcal R(c)}{\log(tc)}
 &=\left(-\frac{\mathcal C(t)}{2L}\mathcal R(c)\right)w(tc),\\
 \frac{\mathcal A(a)\mathcal B(s)}{\log(as)}
 &=\left(\mathcal A(a)\frac{\mathcal B(s)}{2L}\right)w(as).
\end{aligned}
\end{equation}
Here \(\log c/(2L)\) is unary, and
\(\|w\|_\infty\le2L/(L-\log4)=O(1)\), and
\[
 \operatorname{Var}_{I_H}(w)
 =\frac{2L}{L-\log4}-\frac{2L}{L+\log(1-\varepsilon)}=O(1).
\]
Thus uniform product-prefix estimates restore the exact weights by
Stieltjes summation, including the minus sign of \(T_2\).
Use multiplicative cells of relative width \(L^{-D_{\mathrm m}}\).
The original-source boundary strips satisfy
\begin{equation}\label{eq:ra-strip}
 E_{\mathrm{strip}}\ll
       \bigl(NL^{-D_{\mathrm m}/2}+N^{1/2}\bigr)L^{15/2}.
\end{equation}
Each unnormalized source multiplicity is at most \(3d(e)\).
Cauchy and \(d(n)^4\le d_{16}(n)\) give \eqref{eq:ra-strip}.
Sum the boundary strips before the modulus expansion.
For a final saving \(A\), fix
\begin{equation}\label{eq:ra-precision-order}
 D_{\mathrm m}=2A+40,\qquad
 T_{\mathrm{save}}>A+C_{\mathrm{tensor}}+2D_{\mathrm m}+40.
\end{equation}
Then fix every Siegel--Walfisz and half-level precision.
Choose the common threshold last.

\begin{lemma}[Product-region reduction for the central rows]
\label{lem:ra-product-region}
Fix \(A>0\) and choose \(D_{\mathrm m}\) and \(T_{\mathrm{save}}\) as in
\eqref{eq:ra-precision-order}. Take one of the four central rows of
Table~\ref{tab:ra-rows}.
Keep its unary cuts, its normalized arrays from
\eqref{eq:ra-restoration} and the coefficient \(\Gamma\). For
\(N/4<u\le(1-\varepsilon)N\) let \(\mathcal S(u)\) be its normalized
discrepancy sum over \(N/4<nm\le u\), without the factor \(w\). Then
\begin{equation}\label{eq:ra-central-rows}
 \sup_{N/4<u\le(1-\varepsilon)N}|\mathcal S(u)|\ll_A NL^{-A},
\end{equation}
and the same bound holds for the sum over \(nm\in I_H\) with the factor
\(w(nm)\).
\end{lemma}
\begin{proof}
\emph{Selector.} Let \(q=pd\) be supported, with the primes of \(d\) in
decreasing order \(\ell_1>\cdots>\ell_k\). As in the proof of
Lemma~\ref{lem:ra-selector}, the cube gates give
\(p\ell_1\cdots\ell_{i-1}\ell_i^2<Q\) for every \(i\), hence
\[
 \frac{\ell_i}{\ell_{i+1}\cdots\ell_k}<\frac Qq\qquad(1\le i\le k).
\]
Consecutive divisors of a squarefree \(d\) have ratio at most
\(\max_i\ell_i/(\ell_{i+1}\cdots\ell_k)\). Write \(d=\ell_1M\).
The divisors of \(d\) form the two chains \(\mathcal D(M)\) and
\(\ell_1\mathcal D(M)\). If \(\ell_1>M\), the only new gap is
\(\ell_1/M\). If \(\ell_1\le M\), the ranges \([1,M]\) and
\([\ell_1,\ell_1M]\) overlap. A gap of the union below \(M\) lies inside a
gap of \(\mathcal D(M)\).
One above \(M\) lies inside a gap of
\(\ell_1\mathcal D(M)\). Induction on \(k\) gives the claim. Since
\(d>Y\), the largest divisor \(s\le Y\) of \(d\) has a successor
\(s'>Y\) with \(s'/s<Q/q\).
This proves \eqref{eq:ra-central-selector}.
The pair \((r,s)\) is determined by \(q\), and \(q=rs\).
Put \(F(r,s)=\Gamma(rs)\) on selected pairs and \(F(r,s)=0\) otherwise.
This coefficient is
one-bounded and supported on squarefree \(rs\) with \((rs,2N)=1\). Its
dyadic shells satisfy \(Q_1\ll N^{v_{\mathrm{RA}}-y}\) and
\(Q_2\ll N^{y}\), where \(v_{\mathrm{RA}}-y=909/2000\) and
\(y=341/5000\).
There are \(O(L^2)\) shell pairs.

\emph{Parents.} In the first central \(T_2\) row the tested variable and
array are \(n=t\) and
\(\mathcal C(t)/(2L)\).
In the second they are \(n=c\) and
\(\mathcal R(c)\). In the central \(T_{\mathrm{II}}\) rows they are
\(n=s\) and \(\mathcal B(s)/(2L)\) in the first.
In the second they are \(n=a\) and \(\mathcal A(a)\). The other variable carries the remaining one-bounded array,
with the sign of \(T_2\). On \(I_H\), every relevant pair of dyadic parents
has \(H_{\mathrm c}M_{\mathrm c}\asymp N\) and
\(N^{.50}\ll H_{\mathrm c}\ll N^{.54}\) with fixed constants.
There are
\(O(L^2)\) such pairs. Lemma~\ref{lem:ra-sw} gives \eqref{eq:ra-full-sw}
for the tested arrays at these parents.

\emph{Restriction.} Let \(a_n\) be a tested array on its parent and \(J\)
an interval.
On the same parent, \(a_n\mathbf1_J(n)\) satisfies
\(\Delta_{a\mathbf1_J}(I;k,b;h)=\Delta_a(I\cap J;k,b;h)\) for every
interval \(I\). It therefore satisfies \eqref{eq:ra-full-sw} with the same
constants, measured at the original parent. The other array may be
restricted to any interval.

\emph{Mesh.} Put \(\delta=L^{-D_{\mathrm m}}\). Divide each parent into
\(O(\delta^{-1})\) intervals with endpoint ratio at most \(1+\delta\).
All unary cut points of the row are endpoints. For fixed \(u\), call
a cell complete if it lies inside \(\{N/4<nm\le u\}\). Any other cell
meeting this region has all its products within a factor \((1+\delta)^2\)
of \(N/4\) or of \(u\). The part of the region it covers lies in the set
\(\mathcal E_u\) of integers \(e\in(N/4,(1-\varepsilon)N]\) with \(|e-N/4|\le3\delta N\) or
\(|e-u|\le3\delta N\).
Here \(\#\mathcal E_u\ll\delta N+1\). For a complete
cell apply Lemma~\ref{hyp:ra-central} with \(x=N\), saving \(T_{\mathrm{save}}\), the
restricted arrays, \(b=N\bmod r\) and \(a=N\bmod rs\). Both classes are
reduced since \((rs,N)=1\), and \(a\equiv b\pmod r\). Thus
\[
 \Bigl|\sum_{r,s}F(r,s)\mathcal D_{\mathrm{cell}}(rs,N)\Bigr|
 \le\sum_{r\sim Q_1}\max_{(b,r)=1}
   \sum_{\substack{s\sim Q_2\\\mu(rs)^2=1,\ (rs,N)=1}}
   \max_{\substack{a\bmod rs\\a\equiv b\ (r)\\(a,rs)=1}}
   |\mathcal D_{\mathrm{cell}}(rs,a)|
 \ll_{T_{\mathrm{save}}} NL^{-T_{\mathrm{save}}}.
\]
There are \(O(L^{2D_{\mathrm m}+4})\) choices of cell and shell pair, so the
 complete cells contribute \(O(NL^{-T_{\mathrm{save}}+2D_{\mathrm m}+4})\).

\emph{Strips.} Sum the incomplete cells together before expanding
the modulus coefficient. For fixed \(e\), the factorizations
\(e=nm\) in the row carry total normalized weight at most \(d(e)\).
Also \eqref{eq:ra-proper-power-bound} bounds the modulus sum by
\(d(N-e)+O(L)\). Their contribution is therefore
\[
 \ll\sum_{e\in\mathcal E_u}d(e)\{d(N-e)+L\}
 \ll(\#\mathcal E_u)^{1/2}N^{1/2}L^{15/2}
 \ll\bigl(N\delta^{1/2}+N^{1/2}\bigr)L^{15/2},
\]
by Cauchy's inequality twice, \(d(n)^4\le d_{16}(n)\) and
\(\sum_{n\le N}d_{16}(n)\ll NL^{15}\). This is \eqref{eq:ra-strip}.

\emph{Conclusion.} With \(D_{\mathrm m}=2A+40\) and
\(T_{\mathrm{save}}>A+C_{\mathrm{tensor}}+2D_{\mathrm m}+40\), the two contributions are
\(O(NL^{-A-C_{\mathrm{tensor}}-36})\) and \(O(NL^{-A-25/2}+N^{1/2}L^{15/2})\), uniformly
in \(u\). This proves \eqref{eq:ra-central-rows}. Since
\(\|w\|_\infty+\operatorname{Var}_{I_H}(w)=O(1)\), Stieltjes summation
against \(\mathcal S\) gives the weighted bound, with the sign of \(T_2\)
kept.
\end{proof}


\subsection{The main coefficient}

For primes \(\ell\nmid2N\), set \(g_N(\ell)=1/(\ell-1)\), and put
\(g_N(\ell)=0\) for \(\ell\mid2N\). Extend \(g_N\) multiplicatively to
squarefree integers, with \(g_N(\ell^j)=0\) for \(j\ge2\). For primes
\(e\equiv N\pmod p\) in the fixed reduced class modulo \(p\), and primes
\(\ell\nmid2Np\), the condition \(\ell\mid(N-e)/p\) has
density \(1/\varphi(\ell)=g_N(\ell)\). Put
\begin{equation}\label{eq:ra-euler-density}
 V_N(z_s)=\prod_{\substack{2<\ell<z_s\\\ell\nmid N}}
       \left(1-\frac1{\ell-1}\right)
 =\frac{2e^{-\gamma}C(N)}{\alpha L}(1+o(1)),
\end{equation}
with \(C(N)\) as in Paper~II, (1.2). For a marked prime \(p\in N^K\) and a
level \(R\), put \(D=R/p\) and \(s_p=\log D/(\alpha L)\).
Let \(P=P_N(z_s)\) be the product of the primes \(\ell<z_s\) with
\(\ell\nmid2N\).

\begin{lemma}[Upper-sieve application]\label{lem:ra-upper-sieve}
\begin{enumerate}
 \item[(a)] The comparison \eqref{eq:ra-positive-comparison} holds.
 \item[(b)] There is an absolute constant \(K_{\mathrm{dim}}\) such that,
 for every even \(N\),
 \begin{equation}\label{eq:ra-product-condition}
  \prod_{w\le\ell<z}(1-g_N(\ell))^{-1}
  \le\frac{\log z}{\log w}\left(1+\frac{K_{\mathrm{dim}}}{\log w}\right)
  \quad(z>w\ge2).
 \end{equation}
 \item[(c)] For the three pairs
 \((R,K)\in\{(R_L,J),(R_{\mathrm{tail}},J_{\mathrm{tail}}),(Q,J_t)\}\),
 uniformly in the mark and in even \(N\),
 \begin{equation}\label{eq:ra-finite-density}
  \sum_{d\mid P}\frac{\mu_d^+(R/p)}{\varphi(d)}
  \le V_N(z_s)\{F(s_p)+O(L^{-1/3})\},
 \end{equation}
 and, with \(R=N^rL^{-B_r}\),
 \begin{equation}\label{eq:ra-mark-density}
  \sum_q\frac{c_{R,K}(q)}{\varphi(q)}
  \le\frac{2e^{-\gamma}C(N)}{\alpha L}
  \left\{\int_K F\left(\frac{r-t}{\alpha}\right)\frac{dt}{t}+o(1)\right\}.
 \end{equation}
\end{enumerate}
\end{lemma}
\begin{proof}
We use Iwaniec's upper bound for the linear sieve
\cite{Iwaniec1980Remainder}. Let \(\mathcal A\) be a finite sequence and
\(D>1\) a level.
Let \(P(z)\) be the product of the sifting primes below \(z\). The
upper weight \(\mu_d^+\) of \cite[p.~313]{Iwaniec1980Remainder} is
\((-1)^k\) when \(d=\ell_1\cdots\ell_k\mid P(z)\),
\(\ell_1>\cdots>\ell_k\), and
\(\ell_1\cdots\ell_{2l}\ell_{2l+1}^3<D\) for every
\(0\le l\le(k-1)/2\).
Otherwise it is \(0\). With \(j=2l+1\) these are the
gates of \eqref{eq:ra-rosser-weight}, so \eqref{eq:ra-rosser-weight} is
this weight on \(d\mid P\). By the definition of \(W^+\)
\cite[p.~310]{Iwaniec1980Remainder} and Lemma~2 there
\cite[p.~313]{Iwaniec1980Remainder},
\[
 \sum_{d\mid P(z)}\mu_d^+|\mathcal A_d|=W^+\ge S(\mathcal A,P,z)
\]
for every \(D>1\), since \(W^+\) is \(S(\mathcal A,P,z)\) plus nonnegative
sifted counts. If the density satisfies \eqref{eq:ra-product-condition}
with a constant \(K\) and \(z\le D\), Lemma~3, (12), of
\cite[p.~313]{Iwaniec1980Remainder} gives, with \(s=\log D/\log z\),
\[
 \sum_{d\mid P(z)}\mu_d^+g(d)
 \le V(z)\{F(s)+O(e^{\sqrt K-s}(\log D)^{-1/3})\}.
\]

(a) Let \(e\) be a prime in one of the pieces \eqref{eq:ra-three-pieces}
and \(p\) a mark with \(pm=N-e\). If \(P^-(m)\ge z_s\), then
\((m,P)=1\). The inequality \(W^+\ge S\) for the one-term sequence
\(\{m\}\) at level \(D=R/p\) gives
\begin{equation}\label{eq:ra-upper-weight-positivity}
 \sum_{\substack{d\mid P\\d\mid m}}\mu_d^+(D)
   \ge\mathbf1_{(m,P)=1}\ge\mathbf1_{P^-(m)\ge z_s}.
\end{equation}
This holds for every \(D>1\), in particular when \(s_p=383/200<2\).
Omitting \((m,Np)=1\) in \eqref{eq:ra-original-s6} only enlarges the
count. The condition \(d\mid m\) is \(e\equiv N\pmod{pd}\), and
\(c_{R,K}(pd)=\mu_d^+(R/p)\) by \eqref{eq:ra-weight-equality}. The pieces
\eqref{eq:ra-three-pieces} partition \([2,(1-\varepsilon)N)\times J\). The
point \(e=N/4\) lies in \(I_L\) and \(t=1/3-\eta_{\mathrm m}\) lies in
\(J_{\mathrm{tail}}\). Summing \eqref{eq:ra-upper-weight-positivity} over
\(e\) and \(p\), each mark with its multiplicity, gives
\(S_6\le A_{\mathrm{mix}}\) by \eqref{eq:ra-mix-definition}.

(b) For odd \(\ell\),
\[
 \Bigl(1-\frac1{\ell-1}\Bigr)^{-1}
 =\Bigl(1-\frac1\ell\Bigr)^{-1}\Bigl(1-\frac1{(\ell-1)^2}\Bigr)^{-1}.
\]
The factors with \(\ell\mid2N\) equal one.
Omitting them only lowers the product. Mertens' theorem, with its uniform error term, bounds the
product of the first factors over \(w\le\ell<z\) by
\((\log z/\log w)(1+K_1/\log w)\). The product of the second factors
over \(\ell\ge\max(w,3)\) is at most \(1+K_2/\log w\). The constants
\(K_1,K_2\) are absolute. This gives \eqref{eq:ra-product-condition} with one
\(K_{\mathrm{dim}}\), independent of \(N\) and of \(\omega(N)\).

(c) The limiting ranges of \(s_p\) for the three pairs are
\[
 \Bigl(\frac{383}{200},\frac{12639}{4480}\Bigr],\qquad
 \Bigl(\frac{383}{200},\frac{1916149}{1000000}\Bigr],\qquad
 \Bigl(\frac{544243}{250000},\frac{43148397}{14000000}\Bigr],
\]
all inside \([383/200,3.083)\subset(1,\infty)\). The factors \(L^{-B_r}\)
move \(s_p\) by \(O(\log L/L)\). Hence \(z_s\le D\) and \(\log D\asymp L\)
for \(N\ge N_0(B_r)\).
Then (12) gives \eqref{eq:ra-finite-density}, with
an error \(O(L^{-1/3})\) uniform in the mark, the pair and \(N\). For
\eqref{eq:ra-euler-density}, use the factorization in (b) for odd
\(\ell<z_s\) and restore the factors with \(\ell\mid N\). The primes
\(\ell\mid N\) with \(\ell\ge z_s\) number at most \(1/\alpha\) and change
the product by a factor \(1+O(1/z_s)\). By \eqref{eq:ra-weight-equality},
\[
 \sum_q\frac{c_{R,K}(q)}{\varphi(q)}
 =\sum_{\substack{p\in N^K\\p\nmid N}}\frac1{p-1}
   \sum_{d\mid P}\frac{\mu_d^+(R/p)}{\varphi(d)}.
\]
The outer weights are positive, so \eqref{eq:ra-finite-density} may be
summed termwise. Since \(F\) is Lipschitz on the compact range above, the
prime number theorem and partial summation give \eqref{eq:ra-mark-density}.
Marks dividing \(N\) number at most \(1/t_0\), and endpoint atoms and the
logarithmic buffers contribute \(o(1)\).
\end{proof}
The delayed branch \(s_p>3\) of \(F\) occurs only for \(c_{\mathrm{trim}}\),
when \(t<t_{\mathrm{RA}}\), where \(t_{\mathrm{RA}}=v_{\mathrm{RA}}-3\alpha\).


Define
\begin{equation}\label{eq:ra-j-functional}
 J_K(r)=\frac1{2e^\gamma\alpha}
          \int_K F\left(\frac{r-t}{\alpha}\right)\frac{dt}{t}.
\end{equation}
The rational endpoint \(\overline J_6\) is as in Paper II
\cite[(A.13)]{PaperII}:
\begin{equation}\label{eq:ra-j6bar-comparison}
\begin{gathered}
 \overline J_6=
 \frac{13148581959157321533513247276782285310758770468886237619585}
      {10^{58}},\\
 J_J(1/2)=2\log(110/57)\le\overline J_6.
\end{gathered}
\end{equation}
The inequality is certified by outward interval evaluation.
No central approximation is rounded.
Since \(W=4(1-\varepsilon)C(N)N/L^2\), the interval masses give
\begin{equation}\label{eq:ra-main-normalized}
\begin{aligned}
 \frac{M_{\mathrm{mix}}}{W}
 &\le w_LJ_J(1/2)
       +w_H\{J_{J_{\mathrm{tail}}}(1/2)+J_{J_t}(v_{\mathrm{RA}})\}+o(1),\\
 w_L&=\frac{2500}{9999},\qquad w_H=\frac{7499}{9999}.
\end{aligned}
\end{equation}
Here \(w_L+w_H=1\).
No mark multiplicity or singular-series factor is suppressed.

On \(J_{\mathrm{tail}}\), the explicit upper-sieve formula gives
\begin{equation}\label{eq:ra-trim-debit}
\begin{aligned}
 0&\le w_H\{J_{J_{\mathrm{tail}}}(1/2)-J_{J_{\mathrm{tail}}}(v_{\mathrm{RA}})\}\\
 &=w_H\int_{1/3-\eta_{\mathrm m}}^{1/3}
       \frac{v_{\mathrm{RA}}-1/2}{t(1/2-t)(v_{\mathrm{RA}}-t)}\,dt\\
 &\le\frac{w_H\eta_{\mathrm m}(v_{\mathrm{RA}}-1/2)}
              {(1/3)(1/2-1/3)(v_{\mathrm{RA}}-1/3)}\\
 &=:d_{\mathrm m}=\frac{5106819}{31557955000}.
\end{aligned}
\end{equation}
The last integrand is increasing on this interval.
Put
\begin{equation}\label{eq:ra-exact-gain}
 \mathfrak g_{\mathrm{RA}}
   =w_H\{J_J(1/2)-J_J(v_{\mathrm{RA}})\}-d_{\mathrm m}.
\end{equation}
Combining \eqref{eq:ra-main-normalized} with \eqref{eq:ra-j6bar-comparison} gives
\begin{equation}\label{eq:ra-b11-integral}
 \frac{M_{\mathrm{mix}}}{W}
       \le\overline J_6-\mathfrak g_{\mathrm{RA}}+o(1).
\end{equation}

The remaining integral is elementary.
The linear-sieve functions $F$ and $f$ are as in
\cite[(1.8)--(1.9), pp.~172--173]{Iwaniec1980Rosser}.
Set \(t_{\mathrm{RA}}=v_{\mathrm{RA}}-3\alpha\). On the ranges here,
\begin{equation}\label{eq:ra-delay-formula}
 \frac{F(s)}{2e^\gamma}
  =\begin{cases}
      1/s,&1\le s\le3,\\[2pt]
      \displaystyle\frac1s\left\{1+
            \int_2^{s-1}\frac{\log(\xi-1)}\xi\,d\xi\right\},&3<s<5.
    \end{cases}
\end{equation}
This is the linear-sieve normalization
\cite[p.~202]{Iwaniec1980Rosser}\cite[(2.23)]{PaperI}.
Write the delayed contribution as
\begin{equation}\label{eq:ra-delayed-integral}
 \mathcal I_{\mathrm{RA}}=
 \int_{t_0}^{t_{\mathrm{RA}}}\frac{dt}{t(v_{\mathrm{RA}}-t)}
 \int_2^{(v_{\mathrm{RA}}-t)/\alpha-1}
                 \frac{\log(\xi-1)}\xi\,d\xi.
\end{equation}
Substitution gives
\begin{equation}\label{eq:ra-gain-elementary}
\begin{aligned}
 \mathfrak g_{\mathrm{RA}}+d_{\mathrm m}
 =w_H\Bigg\{&\int_{t_0}^{1/3}
   \left(\frac1{t(1/2-t)}-\frac1{t(v_{\mathrm{RA}}-t)}\right)dt\\
 &-\mathcal I_{\mathrm{RA}}\Bigg\}.
\end{aligned}
\end{equation}
Swap the order in \eqref{eq:ra-delayed-integral} to obtain
\begin{equation}\label{eq:ra-swapped-integral}
\begin{aligned}
 \mathcal I_{\mathrm{RA}}&=\frac1{v_{\mathrm{RA}}}
 \int_3^{s_{\max}}h(\zeta)\,d\zeta,\qquad
 s_{\max}=\frac{v_{\mathrm{RA}}-t_0}{\alpha},\\
 h(\zeta)&=\frac{\log(\zeta-2)}{\zeta-1}
 \log\frac{(v_{\mathrm{RA}}-\alpha\zeta)(v_{\mathrm{RA}}-t_0)}
               {\alpha\zeta t_0}.
\end{aligned}
\end{equation}
For \(n=20000\), set \(d=(s_{\max}-3)/n\). The composite midpoint sum
has error at most
\[
 \left|\int_3^{s_{\max}}h(\zeta)\,d\zeta
      -d\sum_{j=0}^{n-1}h(3+(j+1/2)d)\right|
 \le\frac{(s_{\max}-3)d^2}{24}\sup|h''|.
\]
Directed interval differentiation on \(400\) covering boxes gives
\(\sup|h''|\le0.720711111\). An independent directed-interval replay
encloses
\(\mathfrak g_{\mathrm{RA}}\) in
\([0.09895983533531812,0.09895983533543704]\).
Outward rounding therefore gives
\begin{equation}\label{eq:ra-gain-enclosure}
 0.0989598353353\le\mathfrak g_{\mathrm{RA}}
                      \le0.0989598353355.
\end{equation}
Thus we may use the certified rational floor
\begin{equation}\label{eq:ra-rational-floor}
 g_{\mathrm{RA,lo}}=\frac{989598353}{10^{10}}=0.0989598353,
 \qquad \mathfrak g_{\mathrm{RA}}\ge g_{\mathrm{RA,lo}}.
\end{equation}
The exact integral \(\mathfrak g_{\mathrm{RA}}\) and its certified floor
\(g_{\mathrm{RA,lo}}\) are distinct.

\begin{proposition}[Sixth-sum source applications]\label{hyp:ra-applications}
On the exact arrays above, the following hold.
\begin{enumerate}
 \item The comparison \eqref{eq:ra-positive-comparison}, and the density
 bound \eqref{eq:ra-mark-density} for the three pairs \((R,K)\).
 \item The three low discrepancy bounds \eqref{eq:ra-three-bv-errors}.
 \item For every fixed \(A>0\), each of the nine rows of
 Table~\ref{tab:ra-rows}, with its exact weight, is \(O_A(NL^{-A})\).
 Consequently the high-coefficient estimate \eqref{eq:ra-ra1} holds.
\end{enumerate}
\end{proposition}
\begin{proof}
Items 1 and 2 are Lemmas~\ref{lem:ra-upper-sieve} and~\ref{lem:ra-n4-bv}.
For item 3 fix \(A\). Choose in this order: \(D_{\mathrm m}=2A+40\);
\(T_{\mathrm{save}}>A+C_{\mathrm{tensor}}+2D_{\mathrm m}+40\); the saving \(A+C_{\mathrm{tensor}}\) in
Lemma~\ref{hyp:ra-l52} and the saving \(T_{\mathrm{save}}\) in Lemma~\ref{hyp:ra-central}.
Then fix the Siegel--Walfisz precisions that these require, the buffers
\eqref{eq:ra-bv-buffer} and \(B\), and the constant \(K_{\mathrm{dim}}\).
Choose the common threshold last. The rows partition \(T_1+T_2+T_{\mathrm{II}}\):
\(T_1\) is all of \(a\le U\). The sums \(T_2\) in \(c\) and \(T_{\mathrm{II}}\)
in \(a>U\) are split at \(N^{.46}\), \(N^{.50}\) and \(N^{.54}\).
Each boundary goes to the interval on its right.

\emph{Direct rows.} These are \(T_1\) and the part of \(T_2\) with
\(c\ge N^{.54}\). Fix the outer variable \(a\), respectively \(t\). Terms
with \((a,q)>1\), respectively \((t,q)>1\), vanish in \(\mathcal K_q\).
Otherwise \(\mathcal K_q(ac)\) is the kernel of \eqref{eq:ra-theorem-r} for
\(c\) in the reduced class \(N\bar a\bmod q\), and similarly for \(t\). The
parent is \(H=N/a\), respectively \(H=N/t\). The free variable runs over
\(H/4<c<(1-\varepsilon)H\), intersected with \(c\ge N^{.54}\) in the
\(T_2\)-row and with any prefix cut.
This interval lies in \((0,2H]\).
Condition (W1) holds.
For \(a\le U\) we have \(N^{2/3}\le H\le N\).
If the right \(T_2\)-row is nonempty, then
\(N^{.54}<N/t=H\le N\), since
\(N^{.54}\le c<(1-\varepsilon)N/t\). Thus Lemma~\ref{hyp:ra-rough} applies
with this \(H\). (Equivalently, split at \(H/2\) into
\((H/4,H/2]\) and \((H/2,(1-\varepsilon)H)\).
Each lies inside a dyadic parent
\((H',2H']\) with \(H'\in\{H/4,H/2\}\ge N^{.54}/4\). This only doubles the
constant.) The weights
\(\log c/\log(ac)\) and \(\mathcal C(t)/\log(tc)\) are bounded and
monotone in \(c\). Partial summation and \eqref{eq:ra-direct-row-total}
with \(T_{\mathrm{save}}>A+2\) bound both rows by \(O(NL^{-A})\).

\emph{Convolution rows.} These are \(T_2\) with \(c<N^{.46}\) (tested
\(\mathcal R(c)\), \(c>N^{1/3}/4\)), \(T_{\mathrm{II}}\) with
\(a<N^{.46}\) (tested \(\mathcal A(a)\), \(a>N^{1/3}\)), and
\(T_{\mathrm{II}}\) with \(a\ge N^{.54}\) (tested \(\mathcal B(s)/(2L)\),
\(N^{1/3}<s<2N^{.46+O(1/L)}\)). Each tested array contains the
strict roughness factor \(\mathcal R(n)=1_{P^-(n)>z}\).
This is explicit
for \(\mathcal R(c)\).
It follows from
\(\mathcal A(a)=\mu(a)\mathcal R(a)\) and from
\(\mathcal B(s)=\mathcal R(s)\sum_{b\mid s,\ b>V}\Lambda(b)\).
Lemma~\ref{lem:ra-sw} gives \eqref{eq:ra-full-sw} on these original
parents. Write
\(\sum_q\Gamma(q)\Phi(q)=\sum_{r,s}F_H(r,s)\Phi(rs)\) as in
\eqref{eq:ra-global-tensor}. Lemma~\ref{lem:ra-selector} places the
joint support in \eqref{eq:ra-l52-caps}. Discard a dyadic shell pair
\(r\sim U_1\), \(s\sim V_1\) if \(F_H\) vanishes identically there.
Otherwise a supported pair in it satisfies
\(U_1V_1<rs<Q\), \(U_1<D_1\), and \(V_1<D_2\).
Hence this occupied
shell satisfies all scale conditions of Lemma~\ref{hyp:ra-l52}.
Lemma~\ref{lem:ra-tensor} expands \(F_H\) on the whole Cartesian shell
into rank-one terms with \(\sum_\nu|c_\nu|\ll L^{79}\).
The unary caps and coprimality are kept. Each rank-one term is estimated on that
whole shell. Its factors need not preserve the joint condition
\(rs<Q\) or the condition \((r,s)=1\), since
Lemma~\ref{hyp:ra-l52} imposes neither. Apply Lemma~\ref{hyp:ra-l52} with saving \(A+C_{\mathrm{tensor}}\) to each
term on each occupied shell pair and each product prefix
\(I=(N/4,u]\). Since
\(C_{\mathrm{tensor}}\ge83\), the total is \(O(NL^{-A})\), uniformly in \(u\). The factor
\(w\) is then restored by Stieltjes summation, as in
\eqref{eq:ra-restoration}.

\emph{Central rows.} The four central rows are bounded by
Lemma~\ref{lem:ra-product-region}.

\emph{Proper powers.} By \eqref{eq:ra-b10} and
\eqref{eq:ra-proper-power-bound},
\(P_{I_H}(\Gamma)=T_1+T_2+T_{\mathrm{II}}+O(N^{1/2+\nu}L^2)\). This
proves \eqref{eq:ra-ra1}.
\end{proof}

\begin{proposition}[Sixth-sum estimate]\label{prop:ra-b11}
Fix \(\sigma\) in the range stated above before taking \(N\to\infty\).
Then, along even integers,
\begin{equation}\label{eq:ra-b11}
 \frac{S_6}{W}\le\overline J_6-\mathfrak g_{\mathrm{RA}}+o(1).
\end{equation}
In particular, the certified floor in \eqref{eq:ra-rational-floor} gives
\begin{equation}\label{eq:ra-b11-floor}
 \frac{S_6}{W}\le\overline J_6-g_{\mathrm{RA,lo}}
                         +E_{\mathrm{RA}}(N;\sigma),
 \qquad E_{\mathrm{RA}}(N;\sigma)\ge0,\qquad
 E_{\mathrm{RA}}(N;\sigma)\longrightarrow0.
\end{equation}
\end{proposition}

\begin{proof}
By Proposition~\ref{hyp:ra-applications}(1), \(S_6\le A_{\mathrm{mix}}\).
Use \eqref{eq:ra-b9-f1}, obtained from \eqref{eq:ra-b9} and
\eqref{eq:ra-cancellation}. Fix \(A>2\), so that \(NL^{-A}=o(W)\). The
three low discrepancies are \(O(NL^{-A-10})\) by
Proposition~\ref{hyp:ra-applications}(2).
The term
\eqref{eq:ra-high-conversion} is \(O_A(NL^{-A})\).
Also \(P_{I_H}(\Gamma)=O_A(NL^{-A})\) by
Proposition~\ref{hyp:ra-applications}(3). Hence
\(S_6/W\le M_{\mathrm{mix}}/W+o(1)\), and \eqref{eq:ra-b11-integral} gives
\eqref{eq:ra-b11}. Enlarge the vanishing error to make it nonnegative. The
certified floor \eqref{eq:ra-rational-floor} gives \eqref{eq:ra-b11-floor}.
\end{proof}

The precisions are fixed in the order \eqref{eq:ra-precision-order}.
All constants precede the common threshold.
Their dependence through prime Siegel--Walfisz may be ineffective.
The final section uses only \eqref{eq:ra-b11-floor}.
It uses the certified rational floor in \eqref{eq:bill-components}.
Its gain is independent of the common coefficient used in the
balanced-semiprime and switched comparisons.
